% ============================================================================
%  F Sossan HES-SO VS, template for written exams.
%  For the answer key, add the [solutions] option below, or build both with
%  ./make-build.sh.  "Below" is the \usepackage line: write [french,solutions]
%  in both branches, and every \sol{} or solution environment is then printed
%  instead of being dropped.
% ============================================================================
\documentclass[11pt,a4paper]{article}

% Loads the package from ./examsheet when it sits next to this file (Overleaf,
% or a copy you were sent), and from the TeX tree when it was installed with
% ./install.sh.  Options: french|english, solutions, keepspace, nobabel.
\IfFileExists{examsheet/examsheet.sty}
  {\usepackage[french]{examsheet/examsheet}}
  {\usepackage[french]{examsheet}}

% ---------------------------------------------------------------------------
%  Exam identity
% ---------------------------------------------------------------------------
\examtitle{Course code -- Final written test}
\examsubtitle{AY 2025--2026}
\examshorttitle{Final test}          % what goes in the running head
\examdate{\today}
\examduration{2 h}

\examinstructions{%
  Écrivez votre prénom et votre nom en haut de chaque feuille. Seules ces
  feuilles doivent être remises. Écrivez vos réponses finales dans les espaces
  prévus. Utilisez du brouillon pour les esquisses et les calculs
  intermédiaires, qui ne doivent pas être remis. Tout écart sera pénalisé.%
}
\exammaterial{%
  Seule une feuille de résumé recto-verso manuscrite est autorisée. Téléphone,
  ordinateur portable, tablette, etc.\ sont interdits. Toute communication est
  interdite.%
}

\begin{document}
\makeexamtitle

% Optional: a marking table listing every question, its points and a box for
% the score.  It fills itself in from the questions below, on the second pass.
% \markingtable

% ===========================================================================
%  QUESTIONS -- replace everything below with your own
%
%  \question{Title}                  points are summed from the subquestions
%  \question[2]{Title}               ... or given here if there are none
%  \subquestion[1.5]{Text}           lettered (a), (b), (c), ...
%
%  Answer space: \shortanswer  \mediumanswer  \longanswer  \verylonganswer
%                \answergrid{n}  \answerlines{n}  \answerblank{5cm}
%
%  \begin{solution} ... \end{solution}  and  \sol{...} are printed only when
%  the package is loaded with [solutions].
% ===========================================================================

\question{A question with subquestions}

\subquestion[1.5]{First thing you ask. Its points are counted into the
question heading and into the exam total automatically.}
\mediumanswer

\subquestion[1]{Second thing you ask.}
\shortanswer
\begin{solution}
  What you expect to read, shown only in the answer key.
\end{solution}

\question[2]{A question with no subquestions}

Ask it here, then leave room for the answer.
\longanswer

\question{True or false}

\begin{truefalse}
  \tfitem[0.5]{A statement that is true.}
  \tfitem[0.5]{A statement that is false.}
\end{truefalse}

\end{document}
